\subsubsection{利用者中心設計}

\term{利用者中心設計}{User-Centered Design} は、利用者とその作業、目的、文脈を深く理解したうえで設計する方法である。単に画面をきれいにすることではなく、利用者が何を達成したいか、どのような環境で使うか、どのような誤りを起こしやすいかを考える。

利用者中心設計では、利用者要求の収集、作業の流れの整理、プロトタイプやモックアップの作成、利用者評価を行う。使いやすさが重要なシステムでは、利用者中心設計を後から追加するのではなく、設計の初期から組み込む必要がある。
